% Rozdział aktywny po restrukturyzacji DT.

\section{Skoki powierzchniowe i ograniczenia rachunku dystrybucji}
\label{sec:dt-jumps-products}
\providecommand{\atlasref}[2]{#2}
Rozwijamy wzór skoku z rozdziału \ref{sec:pochodne20}
na przypadek hiperpowierzchni. Por.
\atlasref{PDE-07::thm:pde7-jump}{RRC-07: wzór skoku}.
\subsection{Pochodne przez hiperpowierzchnię}
\begin{twierdzenie}[Wzór skoku przez hiperpowierzchnię]
\label{thm:dt-hypersurface-jump}
Niech $d\ge1$, $\K\in\{\R,\C\}$, a $\Omega\subseteq\R^d$
będzie niepustym zbiorem otwartym. Niech
$S\subseteq\Omega$ będzie względnie domkniętą,
zorientowaną, zanurzoną hiperpowierzchnią klasy $C^1$
bez brzegu w $\Omega$. Zakładamy, że w każdym
punkcie $S$ istnieje otoczenie, w którym $S$ jest
wykresem funkcji klasy $C^1$ (po obrocie współrzędnych)
i rozdziela to otoczenie na dwie strony $\Omega_-$
i $\Omega_+$. Nazwy stron są zgodne z orientacją:
jednostkowa normalna $\nu$ jest skierowana
od strony minus do strony plus.
Niech $u\in C^1(\Omega\setminus S;\K)$ i niech lokalnie
przy każdym punkcie $S$ ograniczenia $u$ do obu stron
mają przedłużenia klasy $C^1$ na ich domknięcia.
Oznaczmy ślady na $S$ przez $u_-$ i $u_+$.
Dla $j=1,\ldots,d$ obie strony poniższego wzoru
należą do $\mathcal D'(\Omega;\K)$.
Wtedy, w sensie dystrybucji,
\begin{equation}\label{eq:dt-hypersurface-jump}
 \partial_jT_u=T_{(\partial_ju)_{\mathrm{reg}}}
 +(u_+-u_-)\nu_j\delta_S,
\end{equation}
gdzie $(\partial_ju)_{\mathrm{reg}}$ oznacza klasyczną
pochodną poza $S$, a dla ciągłej funkcji $b:S\to\K$ definiujemy
dystrybucję ważonej miary powierzchniowej przez
\[
 \langle b\delta_S,\varphi\rangle
 =\int_S b(s)\varphi(s)\,\mathrm d\sigma(s),
 \qquad \sigma=\mathcal H^{d-1}\!\restriction_S.
\]
Wartości funkcji $u$ i $(\partial_ju)_{\mathrm{reg}}$ na $S$
nie wpływają na ich dystrybucje regularne.
Zapis $b\delta_S$ oznacza ważoną miarę, nie mnożenie
dystrybucji przez dowolną funkcję ciągłą na $\Omega$.
\end{twierdzenie}
\begin{proof}
W lokalnych współrzędnych wykresowych hiperpowierzchnia
ma miarę Lebesgue'a zero z twierdzenia Fubiniego:
każdy przekrój w kierunku osi wykresu jest co najwyżej
jednopunktowy. Przeliczalne pokrycie takimi otoczeniami
daje $\lambda_d(S)=0$. Na każdym zwartym $K\Subset\Omega$
zbiór $K\cap S$ jest zwarty dzięki względnemu domknięciu
$S$. Można go pokryć skończenie wieloma mniejszymi
otoczeniami wykresowymi o ograniczonej podstawie
i ograniczonej pochodnej funkcji wykresu.
Wzór na miarę powierzchni wykresu
$\sqrt{1+|\nabla h|^2}\,\mathrm d\lambda_{d-1}$
daje $\sigma(K\cap S)<\infty$.
W wymiarze $d=1$ stosujemy zamiast tego miarę
liczącą: $K\cap S$ jest skończonym zbiorem punktów.
Z lokalnych przedłużeń klasy $C^1$ wynika,
że $u$ i jej pochodne klasyczne poza $S$ są
ograniczone na dostatecznie małych otoczeniach
punktów $S$ po obu stronach wykresu. Poza $S$
jest tak z ciągłości. Skończone pokrycie $K$
takimi otoczeniami daje ograniczoność na $K\setminus S$,
a stąd lokalną całkowalność $u$
i $(\partial_ju)_{\mathrm{reg}}$ po nadaniu im
wartości zerowych na $S$.
Jeśli $K\cap S\ne\varnothing$, to
\[
 |\langle b\delta_S,\varphi\rangle|
 \le \sup_{K\cap S}|b|\,\sigma(K\cap S)p_{K,0}(\varphi)
 \qquad(\varphi\in\mathcal D_K(\Omega)).
\]
Supremum jest skończone z ciągłości $b$ na
zwartym $K\cap S$. Gdy przecięcie jest puste,
funkcjonał znika na $\mathcal D_K(\Omega)$. Zatem wyraz powierzchniowy
jest liniowy i spełnia kryterium ciągłości
z Twierdzenia~\ref{thm:distribution-estimate19},
ponieważ $p_{K,0}(\varphi)=p_0(\varphi)$ dla
$\varphi\in\mathcal D_K(\Omega)$.
Ma rząd globalny zero w sensie Definicji~\ref{def:order19}.

Ustalmy $\varphi\in\mathcal D(\Omega)$.
Dla każdego punktu $S$ wybieramy najpierw
dostatecznie mały cylinder wykresowy $C$
o podstawie prostokątnej, ze zwartym domknięciem
w otoczeniu, w którym obowiązują założenia.
W lokalnych współrzędnych zapisujemy
$C=D\times(c_-,c_+)$, gdzie $D$ jest otwartym
prostokątem w $\R^{d-1}$, a wykres $x_d=h(x')$
spełnia $c_-<h(x')<c_+$ na $\overline D$.
Taki wybór jest możliwy dzięki ciągłości $h$
i zmniejszeniu podstawy wokół danego punktu.
Oznaczamy części cylindra po obu stronach
wykresu przez $C_-$ i $C_+$ zgodnie z orientacją.
Najpierw rozważmy test $\varphi$ o nośniku
zawartym w jednym z tak wybranych cylindrów.
Wzór Gaussa na obu częściach cylindra stosujemy
do pola $u\varphi\,e_j$, gdzie $e_j$ jest $j$-tym
wektorem standardowej bazy. Obie części mają
brzeg kawałkami klasy $C^1$, a pole wraz z pierwszymi
pochodnymi przedłuża się ciągle do odpowiedniego brzegu.
Dla funkcji zespolonych stosujemy wzór Gaussa osobno
do części rzeczywistej i urojonej pola.
Obrót współrzędnych nie zmienia wzoru Gaussa
dla pola zapisanego w pierwotnych współrzędnych.
W wymiarze $d=1$ cylinder zastępujemy przedziałem
zawierającym dokładnie jeden punkt $S$ i stosujemy
zwykłe całkowanie przez części na jego dwóch
podprzedziałach, z tymi samymi znakami normalnych.
Na sztucznych ścianach cylindra pole znika,
ponieważ $\varphi$ ma nośnik we wnętrzu.
Normalna zewnętrzna części minus na $S$ to $\nu$,
a części plus to $-\nu$. W konsekwencji
\[
 \begin{aligned}
 -\int_{C_-}u(x)\partial_j\varphi(x)\,\mathrm d\lambda_d(x)
 &=\int_{C_-}(\partial_ju)(x)\varphi(x)\,\mathrm d\lambda_d(x)
   -\int_{S\cap C} u_-\varphi\nu_j\,\mathrm d\sigma,\\
 -\int_{C_+}u(x)\partial_j\varphi(x)\,\mathrm d\lambda_d(x)
 &=\int_{C_+}(\partial_ju)(x)\varphi(x)\,\mathrm d\lambda_d(x)
   +\int_{S\cap C} u_+\varphi\nu_j\,\mathrm d\sigma.
 \end{aligned}
\]
Tutaj całki dotyczą części lokalnego cylindra.
Sumując, otrzymujemy~\eqref{eq:dt-hypersurface-jump} dla takiego
testu. Dla testu o nośniku rozłącznym z $S$ wzór
wynika z klasycznego całkowania przez części.

W ogólnym przypadku pokrywamy zwarty nośnik
$\varphi$ skończoną rodziną uprzednio wybranych
cylindrów wykresowych oraz otoczeń rozłącznych
z $S$. Korzystamy z gładkiego
podziału jedności podporządkowanego temu pokryciu
w otoczeniu nośnika. Dokładniej, wybieramy skończenie
wiele nieujemnych funkcji testowych $\eta_i$
o nośnikach w tych otoczeniach, których suma $A$
jest dodatnia w otoczeniu nośnika $\varphi$
(istnienie wynika z lematu~\ref{lem:cutoffs17}:
dla każdego punktu nośnika wybieramy kulę o domknięciu
zawartym w jednym otoczeniu pokrycia i nieujemną funkcję
odcięcia równą $1$ na mniejszej kuli; zwartość pozwala
wybrać skończenie wiele takich kul).
Funkcje $\varphi_i=\varphi\eta_i/A$ tam, gdzie $A>0$,
przedłużamy przez zero poza ten zbiór.
Są gładkie, gdyż nośnik $\varphi$ jest zawarty
w otwartym zbiorze $\{A>0\}$, oraz mają zwarte
nośniki w odpowiednich otoczeniach.
Mamy $\varphi=\sum_i\varphi_i$. Stosując wzór
lokalny do każdej $\varphi_i$ i sumując skończenie
wiele równości, otrzymujemy tezę dla $\varphi$.
\end{proof}

\begin{wniosek}[Skok strumienia i równanie ciągłości]
Niech $F=(F_1,\ldots,F_d)$ będzie polem rzeczywistym,
którego każda składowa spełnia założenia
twierdzenia~\ref{thm:dt-hypersurface-jump}. Wtedy
\[
 \operatorname{div}T_F
 =T_{(\operatorname{div}F)_{\mathrm{reg}}}
 +[F]\cdot\nu\,\delta_S,\qquad [F]=F_+-F_-.
\]
Równanie $\operatorname{div}T_F=0$ jest równoważne
warunkom $\operatorname{div}F=0$ poza $S$ oraz
$[F]\cdot\nu=0$ na $S$.
\end{wniosek}
\begin{proof}
Sumujemy wzór~\eqref{eq:dt-hypersurface-jump} dla składowych
$F_j$ i pochodnych $\partial_j$.
Jeśli $\operatorname{div}T_F=0$, testy o nośnikach
w $\Omega\setminus S$ dają zerowanie klasycznej
dywergencji prawie wszędzie na tym zbiorze na mocy
stwierdzenia~\ref{prop:regular-injective19}; z ciągłości
jest ona zerowa punktowo. Ponieważ $\lambda_d(S)=0$,
pozostaje $b\delta_S=0$, gdzie $b=[F]\cdot\nu$ jest
ciągłą funkcją rzeczywistą na $S$.
Gdyby $b(s_0)\ne0$, w dostatecznie małym otoczeniu
$s_0$ funkcja $b$ miałaby stały znak i moduł
ograniczony od dołu przez dodatnią stałą.
Wybierając nieujemną funkcję testową o nośniku
w tym otoczeniu, dodatnią w otoczeniu $s_0$,
otrzymalibyśmy niezerową całkę powierzchniową:
każdy niepusty otwarty fragment wykresu ma dodatnią
miarę powierzchniową (dla $d=1$ miarę liczącą).
Sprzeczność dowodzi $b=0$. Odwrotnie, gdy oba
warunki zachodzą, oba składniki wzoru znikają.
Jest to mechanizm stojący za warunkiem
Rankine'a--Hugoniota w teorii praw zachowania.
\end{proof}

\subsection{Funkcje charakterystyczne i pochodne drugiego rzędu}

\begin{wniosek}[Pochodna funkcji charakterystycznej]
\label{cor:dt-characteristic-boundary}
Niech $G\subseteq\R^d$ będzie zbiorem otwartym
o brzegu klasy $C^1$, to znaczy lokalnie brzeg jest
wykresem klasy $C^1$, a $G$ w tym otoczeniu jest
dokładnie jedną z dwóch otwartych stron wykresu.
Oznaczmy przez $n$ jednostkową normalną zewnętrzną do $G$.
Wtedy
\[
 \partial_j T_{\chi_G}=-n_j\delta_{\partial G},
 \qquad j=1,\ldots,d.
\]
W szczególności dla pola testowego
$X\in C_c^\infty(\R^d;\R^d)$
\[
 \int_G \operatorname{div}X(x)\,\mathrm d\lambda_d(x)
 =\int_{\partial G} X(s)\cdot n(s)\,
 \mathrm d\mathcal H^{d-1}(s).
\]
\end{wniosek}
\begin{proof}
Stosujemy Twierdzenie~\ref{thm:dt-hypersurface-jump}
do $S=\partial G$, wybierając stronę minus wewnątrz $G$,
stronę plus na zewnątrz i normalną $\nu=n$.
Ślady $\chi_G$ są odpowiednio $1$ i $0$,
a klasyczne pochodne poza $S$ są zerowe.
Zatem wzór skoku daje pierwszą równość.
Z definicji pochodnej dystrybucyjnej
\[
 -\int_G\partial_j X_j(x)\,\mathrm d\lambda_d(x)
 =-\int_{\partial G}n_j(s)X_j(s)\,
 \mathrm d\mathcal H^{d-1}(s).
\]
Sumując po $j$ i mnożąc przez $-1$, otrzymujemy
drugą równość. Nie wymagamy ograniczoności $G$:
wszystkie całki dotyczą zwartego nośnika pola testowego.
Jest to dystrybucyjna postać wzoru Gaussa,
a nie niezależny dowód twierdzenia Gaussa,
którego użyliśmy w dowodzie wzoru skoku.
\end{proof}

\begin{przyklad}[Kula i miara powierzchniowa]
\label{ex:dt-ball-boundary}
Dla $G=B(a,R)\subseteq\R^d$, $R>0$, normalna
zewnętrzna w punkcie $s\in\partial G$ wynosi
$n(s)=(s-a)/R$, ponieważ jest to jednostkowy
wektor radialny skierowany na zewnątrz kuli.
Z Wniosku~\ref{cor:dt-characteristic-boundary}
\[
 \langle\partial_j T_{\chi_{B(a,R)}},\varphi\rangle
 =-\frac1R\int_{\partial B(a,R)}
 (s_j-a_j)\varphi(s)\,\mathrm d\mathcal H^{d-1}(s).
\]
Dla każdego $j=1,\ldots,d$ ta dystrybucja
nie jest regularna. Istotnie, gdyby
$\partial_jT_{\chi_G}=T_f$ dla
$f\in L^1_{\mathrm{loc}}$, to z zerowania dystrybucji
poza $\partial G$ i Stwierdzenia~\ref{prop:regular-injective19}
mielibyśmy $f=0$ prawie wszędzie poza $\partial G$.
Brzeg ma miarę Lebesgue'a zero, więc $T_f=0$.
Tymczasem na małym otoczeniu punktu $a+Re_j$,
gdzie $n_j>0$, nieujemna funkcja testowa dodatnia
w tym punkcie daje ściśle ujemną wartość powyższego
parowania. Istnienie testu wynika z Lematu~\ref{lem:cutoffs17},
a dodatniość miary fragmentu brzegu z lokalnego wzoru
na miarę wykresu użytego w dowodzie wzoru skoku.
\end{przyklad}

\begin{twierdzenie}[Drugie pochodne i laplasjan przy skoku]
\label{thm:dt-second-interface}
Niech $\Omega,S,u,\nu$ spełniają założenia
Twierdzenia~\ref{thm:dt-hypersurface-jump}.
Załóżmy dodatkowo, że $u$ jest klasy $C^2$ poza $S$,
a lokalne przedłużenia do domknięć obu stron są klasy $C^2$.
Piszemy $[u]=u_+-u_-$ oraz
$[\partial_j u]=(\partial_j u)_+-(\partial_j u)_-$.
Wtedy dla $i,j=1,\ldots,d$ zachodzi
\[
 \partial_i\partial_j T_u
 =T_{(\partial_i\partial_j u)_{\mathrm{reg}}}
 +[\partial_j u]\nu_i\delta_S
 +\partial_i([u]\nu_j\delta_S).
\]
W szczególności
\[
 \Delta T_u
 =T_{(\Delta u)_{\mathrm{reg}}}
 +[\nabla u]\cdot\nu\,\delta_S
 +\operatorname{div}([u]\nu\delta_S).
\]
Ostatni składnik jest określony przez
\[
 \langle\operatorname{div}([u]\nu\delta_S),\varphi\rangle
 =-\int_S [u](s)\nu(s)\cdot\nabla\varphi(s)\,
 \mathrm d\mathcal H^{d-1}(s).
\]
Jeżeli $[u]=0$ na $S$, to
\[
 \Delta T_u=T_{(\Delta u)_{\mathrm{reg}}}
 +[\partial_\nu u]\delta_S,\qquad
 [\partial_\nu u]=[\nabla u]\cdot\nu.
\]
\end{twierdzenie}
\begin{proof}
Wszystkie klasyczne pochodne $u$ rzędów nie większych
niż dwa są lokalnie całkowalne, przez ten sam argument
z lokalnymi przedłużeniami i zwartymi zbiorami co w
Twierdzeniu~\ref{thm:dt-hypersurface-jump}.
Ślady $[u]$ i $[\partial_j u]$ są ciągłe na $S$,
podobnie jak $\nu$. Zatem ich iloczyny z odpowiednimi
składowymi normalnej definiują dystrybucje powierzchniowe
rzędu zerowego przez wzór z tego twierdzenia.
Stosujemy $\partial_i$ do równości
\[
 \partial_j T_u=T_{(\partial_j u)_{\mathrm{reg}}}
 +[u]\nu_j\delta_S.
\]
Do pierwszego składnika stosujemy ponownie wzór skoku,
tym razem dla funkcji $(\partial_j u)_{\mathrm{reg}}$,
która spełnia założenia klasy $C^1$ dzięki założeniu
klasy $C^2$ dla $u$. Otrzymujemy wzór na drugą pochodną.
Sumując go dla $i=j$, dostajemy wzór na laplasjan,
gdzie dywergencja dystrybucji wektorowej oznacza sumę
pochodnych jej składowych. Definicja pochodnej
dystrybucyjnej daje podany wzór na parowanie z testem.
Jeśli $[u]=0$, dystrybucja wektorowa $[u]\nu\delta_S$
jest zerowa, podobnie jak jej dywergencja.
\end{proof}

\begin{przyklad}[Ciągły profil z załamaniem i skok funkcji]
\label{ex:dt-flat-interface}
Niech $S=\{x\in\R^d:x_d=0\}$, $\nu=e_d$,
a $d\ge2$. Dla $u(x)=|x_d|$ ślady $u$ są równe zero,
klasyczny laplasjan poza $S$ jest zerowy, a
$[\nabla u]=2e_d$. Z Twierdzenia~\ref{thm:dt-second-interface}
wynika
\[
 \Delta T_{|x_d|}=2\delta_S,\qquad
 \langle\delta_S,\varphi\rangle
 =\int_{\R^{d-1}}\varphi(x',0)\,\mathrm d\lambda_{d-1}(x').
\]
Dla $v(x)=H(x_d)$ mamy natomiast $[v]=1$,
$[\nabla v]=0$ i klasyczny laplasjan równy zero poza $S$.
Zatem
\[
 \Delta T_v=\partial_d\delta_S,\qquad
 \langle\Delta T_v,\varphi\rangle
 =-\int_{\R^{d-1}}(\partial_d\varphi)(x',0)\,
 \mathrm d\lambda_{d-1}(x').
\]
Pierwszy przykład daje miarę powierzchniową, drugi
jej pochodną w kierunku normalnym. Ta druga dystrybucja
nie jest rzędu zerowego: wybierzmy
$\eta\in\mathcal D(\R^{d-1})$ o całce $1$ i
$\theta\in\mathcal D(\R)$ z $\theta'(0)=1$
(np. $\theta(t)=t\zeta(t)$, gdzie
$\zeta\in\mathcal D(\R;\R)$ jest równa $1$ w otoczeniu $0$).
Funkcję $\eta$ otrzymujemy przez normalizację nieujemnego
niezerowego testu, a $\zeta$ z Lematu~\ref{lem:cutoffs17}.
Dla $\varphi_\varepsilon(x',x_d)
=\eta(x')\varepsilon\theta(x_d/\varepsilon)$,
$0<\varepsilon\le1$, nośniki są zawarte w jednym
zbiorze zwartym, a $p_0(\varphi_\varepsilon)\to0$.
Jednocześnie
$\langle\partial_d\delta_S,\varphi_\varepsilon\rangle=-1$.
To przeczy oszacowaniu rzędu zerowego na tym zbiorze.
Z drugiej strony dla zwartego $K\subseteq\R^d$
i $\varphi\in\mathcal D_K(\R^d)$ zachodzi
\[
 |\langle\partial_d\delta_S,\varphi\rangle|
 \le\lambda_{d-1}(\pi(K))p_1(\varphi),
 \qquad \pi(x',x_d)=x'.
\]
Istotnie, całkowana pochodna znika poza zwartym
rzutem $\pi(K)$ i jej moduł nie przekracza $p_1(\varphi)$.
Miara tego rzutu jest skończona.
Z Definicji~\ref{def:order19} rząd globalny wynosi
więc dokładnie $1$.
\end{przyklad}

\subsection{Przeszkoda Schwartza i kwadrat delty}

Mnożenie dystrybucji przez funkcję gładką jest
określone w definicji \ref{def:multiply20}.
Nie istnieje jednak rozszerzenie tej operacji
na wszystkie dystrybucje z pełnym zestawem
klasycznych własności iloczynu. Poniższe argumenty odróżniają dwa rodzaje przeszkód
algebraicznych od niepowodzenia aproksymacji.
\begin{stwierdzenie}[Przeszkoda w rozszerzeniu iloczynu]
\label{prop:dt-schwartz-obstruction}
Na $\mathcal D'(\R)$ nie istnieje przemienny, łączny
i dwuliniowy iloczyn $\star$ spełniający jednocześnie
następujące warunki:
\begin{enumerate}
\item $T_f\star T_g=T_{fg}$ dla wszystkich
$f,g\in L^\infty_{\mathrm{loc}}(\R)$;
\item $(S\star T)'=S'\star T+S\star T'$ dla wszystkich
$S,T\in\mathcal D'(\R)$.
\end{enumerate}
\end{stwierdzenie}
\begin{proof}
Niech $A=T_H$, gdzie $H$ jest funkcją Heaviside'a.
Ponieważ $H^2=H$ prawie wszędzie, pierwszy warunek
daje $A\star A=A$. Z przykładu~\ref{ex:heaviside20}
mamy $A'=\delta_0$. Różniczkując równość
$A\star A=A$ za pomocą drugiego warunku i korzystając
z przemienności, otrzymujemy
\[
 2(A\star\delta_0)=\delta_0.
\]
Mnożymy tę równość przez $A$. Dwuliniowość i łączność
dają
\[
 A\star\delta_0
 =2A\star(A\star\delta_0)
 =2(A\star A)\star\delta_0
 =2A\star\delta_0.
\]
Stąd $A\star\delta_0=0$, a pierwsza równość
wymusza $\delta_0=0$. Jest to sprzeczne
z działaniem delty na test o wartości $1$ w zerze.
\end{proof}

\begin{stwierdzenie}[Przeszkoda wynikająca z mnożenia przez funkcje gładkie]
\label{prop:dt-smooth-product-obstruction}
Nie istnieje przemienny, łączny, dwuliniowy iloczyn
$\star$ na $\mathcal D'(\R)$ taki, że
\[
 T_g\star T=gT
 \qquad(g\in C^\infty(\R;\K),\ T\in\mathcal D'(\R;\K)),
\]
gdzie prawa strona oznacza mnożenie z
Definicji~\ref{def:multiply20}.
Nie jest tu potrzebna reguła Leibniza.
\end{stwierdzenie}
\begin{proof}
Niech $P=\operatorname{pv}(1/x)$ będzie dystrybucją
wartości głównej z Przykładu~\ref{ex:principal-value19}.
Dla $\varphi\in\mathcal D(\R)$ z definicji mnożenia
i wartości głównej mamy
\[
 \langle xP,\varphi\rangle
 =\lim_{\varepsilon\downarrow0}
 \int_{|x|>\varepsilon}\varphi(x)\,\mathrm d\lambda_1(x)
 =\int_\R\varphi(x)\,\mathrm d\lambda_1(x).
\]
Ostatnia równość wynika z oszacowania modułu całki
po $[-\varepsilon,\varepsilon]$ przez
$2\varepsilon p_0(\varphi)$. Zatem $xP=T_1$.
Jednocześnie $x\delta_0=0$, ponieważ
$\langle\delta_0,x\varphi\rangle=0$.
Z żądanej zgodności z mnożeniem przez funkcje gładkie
$T_1$ jest elementem neutralnym, $T_x\star P=T_1$
oraz $T_x\star\delta_0=0$. Stąd
\[
 \delta_0=T_1\star\delta_0
 =(T_x\star P)\star\delta_0
 =P\star(T_x\star\delta_0)=0.
\]
Ostatni krok wynika z dwuliniowości, a przegrupowanie
z przemienności i łączności. Sprzeczność z
$\delta_0\ne0$ dowodzi tezy.
\end{proof}

Powyższe dwa stwierdzenia dotyczą różnych założeń
o zgodności iloczynu z dotychczasowym rachunkiem:
pierwsze wymaga zachowania iloczynu funkcji lokalnie
ograniczonych i reguły Leibniza, drugie zachowania
mnożenia każdej dystrybucji przez funkcję gładką.
Nie wykluczają one częściowo określonych iloczynów
ani osadzenia dystrybucji w większej algebrze,
w której rezygnuje się z odpowiednich wymagań zgodności.

\begin{przyklad}[Skąd bierze się problem z $\delta^2$?]
\label{ex:dt-delta-square}
Niech $\rho\in\mathcal D(\R;\R)$ będzie nieujemna,
$\supp\rho\subseteq[-1,1]$ i
$\int_\R\rho\,\mathrm d\lambda_1=1$.
Dla $\varepsilon>0$ określamy
$\rho_\varepsilon(x)=\varepsilon^{-1}\rho(x/\varepsilon)$.
Z twierdzenia~\ref{thm:strong-mollification21}
(dla standardowej aproksymacji jedności) wynika
$T_{\rho_\varepsilon}=\delta_0*T_{\rho_\varepsilon}
\to\delta_0$ silnie w $\mathcal D'(\R)$;
równość wynika z Przykładu~\ref{ex:delta-convolution21}.
Bezpośrednio zbieżność słabą-* widać także stąd, że
dla testu $\psi$ mamy
$\langle T_{\rho_\varepsilon},\psi\rangle=
\int\rho(y)\psi(\varepsilon y)\,\mathrm d\lambda_1(y)
\to\psi(0)$ przez jednostajną zbieżność na nośniku $\rho$.
Wybieramy rzeczywisty, nieujemny test $\varphi$
taki, że $\varphi(0)>0$.
Podstawienie $x=\varepsilon y$ daje
\[
 \varepsilon\langle T_{\rho_\varepsilon^2},\varphi\rangle
 =\int_\R \rho(y)^2\varphi(\varepsilon y)\,
 \mathrm d\lambda_1(y).
\]
Ponieważ nośnik $\rho$ jest zwarty,
$\varphi(\varepsilon y)\to\varphi(0)$ jednostajnie
na nim. Po odjęciu liczby
\[
 c=\varphi(0)\int_\R\rho(y)^2\,\mathrm d\lambda_1(y)
\]
moduł różnicy jest ograniczony przez
\[
 \left(\int_\R\rho(y)^2\,\mathrm d\lambda_1(y)\right)
 \sup_{y\in\supp\rho}|\varphi(\varepsilon y)-\varphi(0)|.
\]
Pierwszy czynnik jest skończony, a drugi dąży
do zera. Zatem prawa strona dąży do $c>0$.
Dodatniość wynika z $\varphi(0)>0$ oraz z faktu,
że ciągła funkcja $\rho$ o całce $1$ nie jest zerowa.
Dla dostatecznie małego $\varepsilon$ otrzymujemy
zatem
$\langle T_{\rho_\varepsilon^2},\varphi\rangle
\ge c/(2\varepsilon)\to+\infty$.
Kwadraty tej aproksymacji delty nie mają więc
granicy w $\mathcal D'(\R)$, ponieważ zbieżność
dystrybucyjna wymaga skończonej granicy działania
na każdym ustalonym teście.
Jest to przeszkoda dla definiowania iloczynu
przez zwykłe przejście do granicy gładkich
aproksymacji, odrębna od powyższego argumentu
algebraicznego.
\end{przyklad}
